\SourceLeaf{352}
\begin{center}\large\bfseries [PROJEKTION DES WÜRFELS.]\end{center}

\Needspace{8\baselineskip}\begin{center}[1.]\end{center}

Sind die komplexen Werte der orthografischen Projektionen dreier gleich langer und zueinander senkrechter Geraden
\[
a,\ b,\ c,
\]
so ist
\[
aa+bb+cc=0,
\]
z.\,B.
\[
a=1+8i,\qquad b=8+i,\qquad c=4-4i.
\]
Allgemein kann man setzen
\[
a=pp+qq,\qquad b=2ipq,\qquad c=ipp-iqq.
\]

\Needspace{9\baselineskip}\begin{center}[2.]\end{center}

Die zierlichste Form der Auflösung der allgemeinen Gleichung
\[
pp+qq+rr=0
\]
lautet, wobei $a$ und $b$ beliebige komplexe Zahlen bezeichnen,
\[
p=(a-b)(b-aii)
\]
\[
q=(b-bi)(aii-ai)
\]
\[
r=(bi-a)(ai-b).
\]

\Needspace{9\baselineskip}\begin{center}\bfseries BEMERKUNGEN.\end{center}

Ein Teil dieser beiden Notizen, die in einem Handbuch stehen und aus dem Jahr 1831 stammen, ist bereits in Bd.\,II, S.\,309 abgedruckt.

\indent Von Dezember 1839 bis Ostern 1840 hielt \textsc{Gauss} eine Vorlesung über die \emph{Theorie der imaginären Grössen}, die E.\,\textsc{Heine} hörte und ausarbeitete, wie bereits auf S.\,330 erwähnt. Aus dieser Ausarbeitung stammen die folgenden Erläuterungen zur Projektion des Würfels.
\par\noindent\hfill\textsc{Stäckel}.\par
